\documentclass[10pt,a4paper]{amsart}
\usepackage[margin=19mm]{geometry}
\usepackage{amsmath,amssymb,mathtools}
\usepackage[scr=rsfs,cal=euler]{mathalfa}
\usepackage{xfrac}
\usepackage[unicode,hidelinks,bookmarks=false]{hyperref}
\newcommand{\ie}{i.e.\ }
\newcommand{\cf}{cf.\ }
\newcommand{\resp}{respectively\ }
\newcommand{\Subsection}{\subsection}
\providecommand{\nobreakdash}{\nobreak\hbox{-}}
\setlength{\emergencystretch}{2em}
\multlinegap0pt
\usepackage{xcolor}
\usepackage{amssymb}
\usepackage[shortlabels]{enumitem}
\newtheorem{theorem}{Theorem}
\newtheorem{prop}{Proposition}
\newcommand{\N}{\mathbf{N}}
\newcommand{\Q}{\mathbf{Q}}
\newcommand{\R}{\mathbf{R}}
\newcommand{\w}{\omega}
\title{full title}
\author{Bongiorno Nicolas}
\date{Source: arXiv:2603.13938v1 (CC BY 4.0)}
\begin{document}
\maketitle

\noindent\emph{Source and licence.} full title. Version arXiv:2603.13938v1 (CC BY 4.0), distributed by arXiv under the Creative Commons Attribution 4.0 International licence, http://creativecommons.org/licenses/by/4.0/, as recorded in the arXiv licence metadata for this exact version and shown on the abstract page https://arxiv.org/abs/2603.13938v1. Full licence notice: \texttt{licenses/arxiv-2603.13938-CC-BY-4.0.txt}.\par\smallskip

\noindent\textit{Editorial scope.} Counted unit: the article's prop proposition\_first\_hypothesis\_arithmetic\_function (source0.tex lines 786--798). The article's complete original proof of it is delivered with it (source0.tex lines 800--842, one \texttt{proof} environment of 43 lines). No mathematics is written, repaired or re-derived: the statement and the proof are the article's own lines, in the article's own notation, and no retained line is substituted or re-flowed.  Retained uncounted as the source-local context the proof needs: lines 732--749.  Nothing was omitted from any retained span, so the package carries no omission record.  No retained line contains a citation, so the wrapper carries no bibliography and no bibliography entry.  No retained line contains a figure environment, an image-inclusion command or drawing code, so no image is delivered and the package declares no asset dependency.  Editorial formatting only, in addition: an AMS \texttt{amsart} preamble that restates the article's own declarations and macros used by the retained lines, namely the environments \texttt{theorem}, \texttt{prop} on separate counters, together with the standard packages those lines need.  The retained environments print in this document as Theorem 1, Proposition 1 in order of appearance, whereas the article prints the same items with the numbering of its own full document.  The complete native source of the article as distributed in the arXiv e-print for this exact version is bundled unchanged as \texttt{sources/196360/source0.tex}.
\begin{theorem}\label{theorem_multi_height_value_in_subcone}
Let \(B = (B_1,\ldots,B_{\rho}) \in (\R_{\geqslant 1})^{\rho}\).
Then there exist constants \(C > 0\) and \(\delta > 0\) such that
\[
\left|
\sharp ( T(\Q) )_{h \in \Lambda \cap (-\Lambda)_{B} }
-
\nu(-\Lambda)\, \tau(X)\,
\prod_{i = 1}^{\rho}
B_i^{\langle \w_{X}^{-1}, L_i^* \rangle}
\right|
\leqslant
C \cdot
\frac{\prod_{i = 1}^{\rho}
B_i^{\langle \w_{X}^{-1}, L_i^* \rangle}}
{\min(B_i)^{\delta}}.
\]
\end{theorem}

\begin{prop}\label{proposition_first_hypothesis_arithmetic_function}
    Let $B_1, \dots, B_{\rho} \in \R_{\geqslant 1}$, with sufficiently large $\min\limits_{1 \leqslant i \leqslant {\rho}}(B_i)$ , we have the estimates:
\begin{itemize}
    \item $\sum\limits_{1 \leqslant y_i \leqslant B_i,\ 1 \leqslant i \leqslant \rho} f_{\uparrow}(y) 
= \nu(-\Lambda) \cdot \tau(X) \cdot \prod\limits_{i = 1}^{\rho} B_i^{\langle \w_X^{-1},L_i^* \rangle} 
+ O \!\left( \prod\limits_{i = 1}^{\rho} B_i^{\langle \w_X^{-1},L_i^* \rangle} \cdot
\min\limits_{1 \leqslant i \leqslant {\rho}}(B_i)^{-\Delta} \right) ;$ 
    \item $\sum\limits_{1 \leqslant y_i \leqslant B_i,\ 1 \leqslant i \leqslant \rho} f_{\downarrow}(y) 
= \nu(-\Lambda) \cdot \tau(X) \cdot \prod\limits_{i = 1}^{\rho} B_i^{\langle \w_X^{-1},L_i^* \rangle} 
+ O \!\left( \prod\limits_{i = 1}^{\rho} B_i^{\langle \w_X^{-1},L_i^* \rangle} \cdot
\min\limits_{1 \leqslant i \leqslant {\rho}}(B_i)^{-\Delta} \right) $.
\end{itemize}
\end{prop}

\begin{proof}
First, using the properties of the ceiling function, we obtain for
\(B \in (\N^{*})^{\rho}\) the equality
\[
\sum_{1 \leqslant y_i \leqslant B_i,\ 1 \leqslant i \leqslant \rho}
f_{\uparrow}(y)
=
\sharp T(\Q)_{h \in \Lambda \cap (-\Lambda)_B }.
\]
In this case, the desired estimate follows directly from
Theorem~\ref{theorem_multi_height_value_in_subcone}.
We conclude by observing that we have the bounds
\[
\sum_{1 \leqslant y_i \leqslant \lfloor B_i \rfloor,\ 1 \leqslant i \leqslant \rho}
f_{\uparrow}(y)
\leqslant
\sum_{1 \leqslant y_i \leqslant B_i,\ 1 \leqslant i \leqslant \rho}
f_{\uparrow}(y)
\leqslant
\sum_{1 \leqslant y_i \leqslant \lceil B_i \rceil,\ 1 \leqslant i \leqslant \rho}
f_{\uparrow}(y).
\]

Similarly, in order to estimate
\[
\sum_{1 \leqslant y_i \leqslant B_i,\ 1 \leqslant i \leqslant \rho}
f_{\downarrow}(y),
\]
it suffices to treat the case where \(B_1, \dots, B_{\rho} \in \N^*\).
Moreover, we have the inequalities
\[
\sharp T(\Q)_{h \in\Lambda \cap (-\Lambda)_B}
\leqslant
\sum_{1 \leqslant y_i \leqslant B_i,\ 1 \leqslant i \leqslant \rho}
f_{\downarrow}(y)
\leqslant
\sharp T(\Q)_{h \in \Lambda \cap (-\Lambda)_{B+ \underline{1} }},
\]
where \(\underline{1}\) denotes here the element of \(\R^{\rho}\) whose coordinates in the
canonical basis are all equal to \(1\).
The estimate is therefore again a consequence of
Theorem~\ref{theorem_multi_height_value_in_subcone}.
\end{proof}

\end{document}
